\chapter*{Introduction générale}
\markboth{Introduction générale}{}
\addcontentsline{toc}{chapter}{Introduction générale}

\section*{Contexte}

Dans un marché des services numériques en forte croissance, les entreprises de
type ESN (Entreprise de Services Numériques) font face à un défi commercial
récurrent~: identifier rapidement et avec précision les organisations qui ont un
besoin \textbf{actif et budgété} en prestations IT. Ce besoin peut se manifester
sous différentes formes~: un appel d'offres public, une demande de proposition
(RFP), une recherche de partenaire technologique, ou encore — de façon indirecte
— le recrutement d'un profil IT rare que l'entreprise peine à trouver en CDI.

L'identification manuelle de ces signaux est chronophage, non systématique, et
fortement dépendante de l'intuition et du réseau personnel des commerciaux.
Plusieurs heures de veille hebdomadaire sur LinkedIn, le BOAMP, et les journaux
officiels ne produisent que quelques leads de qualité incertaine.

\section*{Objectifs du projet}

Ce rapport présente \textbf{LeadGen Francophone 360+}, une plateforme de
génération automatisée de leads B2B IT, développée pour Solvinya Group et
ciblant les marchés francophone~: France, Belgique, Luxembourg, Suisse, et
Tunisie.

Les objectifs principaux sont~:
\begin{enumerate}[label=\textbf{\arabic*.}]
  \item \textbf{Collecter} automatiquement les signaux B2B depuis 5 sources
        (LinkedIn Jobs, LinkedIn B2B Posts, BOAMP, TED EU, Malt.fr) ;
  \item \textbf{Qualifier} chaque prospect avec un score 0--100 basé sur des
        règles métier déterministes (pays, technologie, secteur) ;
  \item \textbf{Classifier} le texte des opportunités par secteur et niveau de
        priorité via un modèle NLP (TF-IDF + LinearSVC) ;
  \item \textbf{Générer} des messages de prospection personnalisés prêts à envoyer
        via un LLM (Groq, modèle llama-3.1-8b-instant) ;
  \item \textbf{Automatiser} l'ensemble du pipeline avec n8n et exposer une API
        REST FastAPI.
\end{enumerate}

\section*{Organisation du rapport}

Ce rapport est structuré en quatre chapitres~:

\begin{description}
  \item[\textbf{Chapitre 1}] expose le cadre du projet~: présentation de Solvinya
    Group, problématique détaillée, solution proposée et méthodologie de
    développement adoptée.

  \item[\textbf{Chapitre 2}] présente l'analyse des besoins fonctionnels et
    non fonctionnels, les cas d'utilisation, et l'architecture globale du système.

  \item[\textbf{Chapitre 3}] décrit la conception technique détaillée~:
    architecture des collecteurs, pipeline d'ingestion, formule de scoring,
    module NLP, générateur LLM, schéma de base de données et workflows n8n.

  \item[\textbf{Chapitre 4}] présente la réalisation~: environnement technique,
    implémentation de chaque composant, tests et résultats obtenus (815 prospects
    collectés au 2 mai 2026).
\end{description}

\newpage
